\section{Эксперимент}
\label{sec:experiment}

\porozhdeno{\Claim{E01}{В исходном сравнении медианы покрытия ридов назначениями на уровне семейства составляют \CmpSilvaFamilyReadMedianPct{}\% для SILVA, \CmpGtdbFamilyReadMedianPct{}\% для GTDB и \CmpGgTwoFamilyReadMedianPct{}\% для GG2 (табл.~\ref{tab:comparison}). В объединённом наборе почвенных ASV доли вариантов, получивших назначение на уровне семейства, составляют соответственно \CmpSilvaFamilyAsvPct{}\%, \CmpGtdbFamilyAsvPct{}\% и \CmpGgTwoFamilyAsvPct{}\%. Знаменателем этих долей служит число уникальных ASV, обнаруженных в почвенных образцах; каждый ASV учитывается один раз. Медианы расстояния $D$ между профилями на уровне семейств равны \CmpSilvaGtdbFamilyDistanceMedian{} для SILVA--GTDB, \CmpSilvaGgTwoFamilyDistanceMedian{} для SILVA--GG2 и \CmpGtdbGgTwoFamilyDistanceMedian{} для GTDB--GG2.}}

\begin{table*}[t]
\centering
\caption{\porozhdeno{\Claim{T01}{Исходное сравнение профилей, построенных по всем ридам после DADA2: медианы покрытия назначениями и попарных расстояний $D$ по почвенным образцам для рангов от типа до рода. Отрицательный контроль исключён. Знаменатель покрытия --- общее число ридов образца, включая риды без назначения на соответствующем ранге. При расчёте $D$ учитывается категория \texttt{UNCLASSIFIED}, поэтому расстояния следует интерпретировать совместно с покрытием назначениями.}}}
\label{tab:comparison}
% scope: comparison_all_post_dada2
% Покрытие ридов median (percent); D median; all post-DADA2 reads.
\begin{tabular}{lrrrrrr}
\hline
Ранг & \multicolumn{3}{c}{Покрытие ридов (\%)} & \multicolumn{3}{c}{D} \\
 & SILVA & GTDB & GG2 & SILVA--GTDB & SILVA--GG2 & GTDB--GG2 \\
\hline
Тип & 73.93 & 42.72 & 47.83 & 0.458208 & 0.442640 & 0.030651 \\
Класс & 67.99 & 41.57 & 47.40 & 0.475849 & 0.476744 & 0.048316 \\
Порядок & 65.76 & 40.03 & 45.07 & 0.507603 & 0.511646 & 0.108306 \\
Семейство & 64.96 & 36.20 & 35.68 & 0.554920 & 0.560384 & 0.167251 \\
Род & 47.36 & 21.79 & 22.94 & 0.444442 & 0.455923 & 0.186607 \\
\hline
\end{tabular}

\medskip
\caption{\porozhdeno{\Claim{T02}{Сравнение на уровне семейств после исключения ридов с явной органелльной аннотацией и отбора ридов, отнесённых каждой моделью к Bacteria или Archaea. «Риды» --- общее число ридов почвенного образца после DADA2; $N$ --- число ридов, отобранных по назначениям данной модели; $U$ --- доля ридов категории \texttt{UNCLASSIFIED} среди $N$, выраженная в процентах. При $N=0$ доля не определена и обозначена NA. Терминальная метка, точно совпадающая с \texttt{f\_\_uncultured}, трактуется как отсутствие назначения на уровне семейства. В отличие от табл.~\ref{tab:comparison}, знаменатель определяется после отбора и может различаться между моделями.}}}
\label{tab:family}
% scope: family_prokaryote_selected
% Reads: post-DADA2; N: selected reads; U: UNCLASSIFIED / N (percent); N=0 gives NA.
\begin{tabular}{lrrrrrrr}
\hline
Образец & Риды & \multicolumn{2}{c}{SILVA} & \multicolumn{2}{c}{GTDB} & \multicolumn{2}{c}{GG2} \\
 & & N & U (\%) & N & U (\%) & N & U (\%) \\
\hline
Kr1\_S80 & 12351 & 38 & 100.00 & 10142 & 100.00 & 10114 & 100.00 \\
Kr2\_S81 & 8904 & 0 & NA & 8013 & 100.00 & 8098 & 100.00 \\
Kr3\_S82 & 9985 & 2 & 100.00 & 9147 & 100.00 & 8974 & 100.00 \\
Kr4\_S83 & 12293 & 249 & 98.80 & 11518 & 99.97 & 12013 & 99.98 \\
Kr5\_S84 & 1505 & 0 & NA & 1413 & 100.00 & 1465 & 100.00 \\
LUB2\_S92 & 25216 & 25197 & 8.66 & 25216 & 14.15 & 25214 & 12.44 \\
LUB3\_S93 & 6623 & 6615 & 8.84 & 6623 & 13.69 & 6618 & 17.62 \\
LUB4\_S94 & 16199 & 16182 & 9.58 & 16195 & 15.50 & 16197 & 19.32 \\
LUB5\_S95 & 19090 & 19084 & 13.08 & 19090 & 27.62 & 19086 & 28.65 \\
LUB6\_S96 & 28821 & 28660 & 8.85 & 28821 & 18.99 & 28821 & 18.15 \\
\hline
\end{tabular}

\end{table*}

\porozhdeno{\Claim{E02}{Среди ASV, получивших назначения на уровне рода от обеих сравниваемых моделей, доля буквальных совпадений составляет \CmpSilvaGtdbGenusConditionalAsvPct{}\% для SILVA--GTDB, \CmpSilvaGgTwoGenusConditionalAsvPct{}\% для SILVA--GG2 и \CmpGtdbGgTwoGenusConditionalAsvPct{}\% для GTDB--GG2. Медианы условной доли совпадения по ридам равны соответственно \CmpSilvaGtdbGenusConditionalReadMedianPct{}\%, \CmpSilvaGgTwoGenusConditionalReadMedianPct{}\% и \CmpGtdbGgTwoGenusConditionalReadMedianPct{}\%. Медианы рассчитаны только по образцам, для которых эта доля определена; для пары SILVA--GTDB число таких образцов составляет \CmpSilvaGtdbGenusDefinedSamples{}. Большее совпадение назначений или меньшее расстояние между профилями сами по себе не свидетельствуют о более высокой точности моделей.}}

\porozhdeno{\Claim{E03}{При отборе Bacteria и Archaea по назначениям SILVA для образцов Kr2\_S81 и Kr5\_S84 получено $N=0$, поэтому доли на уровне семейств не определены. При отборе по назначениям GTDB и GG2 для этих образцов $N>0$, однако назначения на уровне семейства отсутствуют (табл.~\ref{tab:family}). Следовательно, необходимо различать отсутствие ридов после отбора и отсутствие назначений на уровне семейства среди отобранных ридов. В образцах LUB назначения на уровне семейства присутствуют у всех трёх моделей. Подробные профили с отдельными легендами приведены на рис.~\ref{fig:overview}--\ref{fig:ggtwo}. Они характеризуют результаты таксономического назначения и сами по себе не позволяют установить экологические причины различий.}}